\section{Knowledge-Guided Approach}
\label{sec:knowledge-guided-approach}

The proposed approach generates a smart-contract-specific threat model by combining contract analysis, domain knowledge, and a structured threat metamodel. The objective is to move beyond generic vulnerability checklists by grounding threat generation in ERC-specific knowledge and by constraining the output through the ERC Threat Metamodel (ETM). As illustrated in Figure~\ref{fig:phase1-approach}, the workflow starts from a smart contract source file or ABI, extracts the contract context, retrieves relevant knowledge, generates an ERC-aware threat model, and validates the resulting model against the ETM.

\subsection{ERC Security Knowledge Base}
\label{subsec:erc-security-kb}

To support the RAG-based ERC threat modeling agent, we constructed a versioned and extensible security knowledge base. The knowledge base combines primary ERC-specific sources with complementary vulnerability, security-control, and risk-assessment knowledge. It includes official ERC specifications collected from the Ethereum ERC repository, smart contract weakness information from SWC and OWASP Smart Contract Security, and auxiliary knowledge from MITRE ATT\&CK, DREAD, CVSS, and NIST SP 800-53. The collected raw artifacts are stored as traceable snapshots with source URLs, repository metadata, timestamps, and SHA-256 hashes. These artifacts are then normalized into structured records covering ERC standards, vulnerabilities, threat patterns, and auxiliary risk/contextualization models. The resulting indexed knowledge base currently contains 600 normalized ERC standards, 207 vulnerability records, 4 auxiliary knowledge records, and 815 retrieval-ready documents. This structure enables the agent to retrieve ERC-specific specifications, weakness descriptions, mitigation guidance, and risk-prioritization knowledge while preserving reproducibility and source traceability.

\begin{table}[t]
\centering
\caption{Structure of the collected ERC security knowledge base}
\label{tab:erc-security-kb}
\begin{tabular}{p{0.18\linewidth} p{0.25\linewidth} p{0.33\linewidth} p{0.16\linewidth}}
\hline
\textbf{KB Layer} & \textbf{Sources} & \textbf{Role in the System} & \textbf{Current Status} \\
\hline
ERC standards & Ethereum ERC repository & Provides official ERC specifications, functions, events, and standard metadata. & 600 normalized ERC records \\
Vulnerability knowledge & SWC Registry; OWASP Smart Contract Security & Provides smart contract weaknesses, descriptions, CWE/SWC/OWASP mappings, and remediation guidance. & 207 vulnerability records \\
Threat patterns & Curated security patterns & Links common smart contract threats to assets, entry points, STRIDE categories, security properties, and test objectives. & Seeded and indexed \\
Auxiliary knowledge & MITRE ATT\&CK; DREAD; CVSS; NIST SP 800-53 & Enriches threats with adversary context, risk scoring, prioritization, and mitigation guidance. & 4 auxiliary records \\
Raw snapshots & GitHub repositories and official web pages & Preserves traceability using source URLs, timestamps, commits, and SHA-256 hashes. & 958 raw files \\
Retrieval index & Normalized KB records & Provides retrieval-ready documents for the RAG agent. & 815 indexed documents \\
\hline
\end{tabular}
\end{table}


\subsection{Contract Context Extraction}
\label{subsec:contract-context-extraction}

The first step consists of extracting the security-relevant context of the target smart contract. The input can be either a Solidity source file or an ABI. From this input, the approach identifies the implemented ERC standard, the exposed functions, their signatures, inheritance information, modifiers, events, and relevant state variables. This context captures the contract-specific surface that may influence the threat model.

ERC standard detection is particularly important because different standards expose different assets, operations, and threat scenarios. For example, an ERC-20 contract is mainly characterized by balances, allowances, and transfer authorization, whereas an ERC-721 contract introduces token ownership, operator approvals, metadata, minting constraints, and safe transfer semantics. The extracted context therefore acts as the bridge between the concrete contract and the standard-specific knowledge stored in the knowledge base.

\subsection{Knowledge Retrieval}
\label{subsec:knowledge-retrieval}

After the contract context has been extracted, the detected ERC standard and contract features are used to retrieve relevant knowledge from the knowledge base. The retrieval process is guided by the standard name, the available functions, and the contract elements identified during analysis. This allows the agent to focus on knowledge that is relevant to the specific contract under analysis rather than using the entire knowledge base indiscriminately.

The retrieved knowledge may include ERC specifications, known vulnerability descriptions, threat patterns, security mappings, mitigations, and security properties. For instance, the retrieval process for an ERC-721 contract prioritizes information related to ownership integrity, approval abuse, operator permissions, unsafe receiver handling, and minting invariants. In contrast, an ERC-20 contract leads to the retrieval of knowledge related to allowance race conditions, transfer authorization, and accounting invariants.

\subsection{ERC-Aware Threat Model Generation}
\label{subsec:erc-aware-generation}

The ERC-aware threat model agent combines three inputs: the extracted contract context, the retrieved knowledge, and the ETM. The ETM defines the conceptual structure of the generated model, including actors, trust levels, assets, entry points, vulnerabilities, threats, countermeasures, security properties, test objectives, and relations between these elements.

The generation process is standard-aware. For widely used ERC standards, the agent relies on explicit standard profiles that capture typical assets and threat categories. For less common or unknown ERC standards, a generic fallback profile is used to generate a conservative threat model from the detected functions and retrieved knowledge. This design makes the approach extensible while ensuring that the generated model remains compatible with the ETM.

The output is an ETM-compliant threat model represented as a structured JSON document. Each threat is connected to the asset it targets, the vulnerability it exploits, the security property it violates, the countermeasure that mitigates it, and the test objective that can be used to validate it. These relations make the threat model suitable for later phases such as attack-tree construction and security-test generation.

\subsection{ETM Validation and Consistency Checking}
\label{subsec:etm-validation}

The generated threat model is validated before it is accepted as an output of the approach. Validation is performed at two complementary levels. First, schema validation ensures that the generated JSON document follows the ETM structure and contains the required sections and fields. Second, consistency checking verifies that the model is semantically connected and internally coherent.

The consistency checks ensure, for example, that each threat targets at least one asset, exploits at least one vulnerability, is mitigated by a countermeasure, violates a security property, and is linked to a test objective. The checker also verifies that relations point to existing model elements, that the generated model matches the detected ERC standard, and that source traceability is available when knowledge items are used. Missing SWC mappings are not treated as automatic errors, since some modern or standard-specific vulnerability classes do not have a precise SWC identifier. In such cases, the model keeps \texttt{SWC: N/A} and relies on other mappings such as CWE, OWASP, or category-level descriptions.

\subsection{Traceability and Reproducibility}
\label{subsec:traceability-reproducibility}

Traceability is a central design goal of the approach. The generated model preserves links to the knowledge sources used during generation, allowing a threat, vulnerability, mitigation, or test objective to be traced back to retrieved knowledge items. This improves explainability and helps distinguish grounded model elements from unsupported claims.

Reproducibility is also preserved through the deterministic generation mode. Given the same contract context and knowledge base version, the agent produces a stable ETM-compliant model. An optional LLM-assisted mode is available for experimental comparison and qualitative enrichment, but the deterministic mode remains the primary generation mode because it is easier to validate, reproduce, and explain.

\section{Implementation Details}
\label{sec:implementation-details}

The Phase~1 prototype implements the knowledge-guided threat modeling workflow as an offline, deterministic, and ETM-validated pipeline. The implementation is organized around five main components: a contract analyzer, a knowledge base retriever, an ERC-aware threat model generator, an ETM validator, and a consistency checker.

The contract analyzer accepts either a Solidity source file or an ABI JSON file. It extracts the contract name, compiler version, functions, signatures, visibility, state mutability, modifiers, events, inheritance declarations, imports, and state variables. It also performs ERC standard detection using multiple signals, including inheritance names, import paths, known function signatures, and interface identifiers when available.

The knowledge base is stored locally and contains normalized records for ERC standards, vulnerabilities, threat patterns, auxiliary security knowledge, and mappings to security taxonomies such as OWASP, SWC, and CWE. The retrieval layer uses the detected ERC standard and extracted contract features to query the knowledge base. Retrieved knowledge items are attached to the generation state and later used as traceable evidence for the produced model.

The threat model generator currently uses a deterministic ERC-aware strategy. Explicit profiles are implemented for common standards such as ERC-20, ERC-721, ERC-1155, and ERC-4626. For other ERC standards, the system uses a generic fallback profile that generates standard-relevant threats based on the detected functions and retrieved knowledge. This allows the prototype to generate an ETM model even when a specific ERC profile has not yet been manually defined.

The generated threat model follows the ETM schema. It contains the following sections: \texttt{ercStandard}, \texttt{smartContract}, \texttt{actors}, \texttt{trustLevels}, \texttt{assets}, \texttt{entryPoints}, \texttt{vulnerabilities}, \texttt{threats}, \texttt{countermeasures}, \texttt{securityProperties}, \texttt{testObjectives}, and \texttt{relations}. The relations connect threats to assets, vulnerabilities, countermeasures, security properties, and test objectives.

After generation, the model is validated in two stages. First, schema validation checks whether the JSON output conforms to the ETM structure. Second, a consistency checker verifies semantic and structural properties, such as whether every threat targets at least one asset, exploits at least one vulnerability, is mitigated by a countermeasure, violates a security property, and is linked to a test objective. The checker also verifies relation integrity, standard matching, traceability warnings, and SWC mapping coverage.

The implementation also includes an optional LLM-assisted generation mode using Groq through an OpenAI-compatible API. However, the deterministic mode remains the primary generation mode because it is reproducible, stable, explainable, and easier to validate. The LLM mode is retained as an experimental variant for future comparison and qualitative enrichment.

For evaluation, a Phase~1 test suite was created with multiple smart contracts implementing OpenZeppelin v5-compatible standards, including ERC-20, ERC-721, ERC-1155, and ERC-4626. For each contract, the pipeline generates a contract context, an ETM-compliant threat model, an agent execution state, and a validation report. This enables systematic assessment of the generated models across different ERC standards.
